\scanpage{192}
\mspage{II.40}
% archive sheet number: 191
% block: top, line 1 | text
Ce dernier donne lieu encore à des hom.

% block: upper third, lines 2-3 | display
\[
\begin{aligned}
\Pi^{\mathrm{an}}&\unsure{\text{---}|\longrightarrow}\Pi^{\mathrm{alg}}\text{---}\Pi\\
&\qquad\deleted{\unsure{\alpha}}
\end{aligned}
\]

% block: upper third, lines 4-5 | display
\[
\begin{aligned}
\text{et}\qquad\Pi&\simeq\Pi^{\mathrm{alg}}/(\Pi^{\mathrm{alg}})^0\\
\mathfrak g=\operatorname{Lie}(\Pi^{\mathrm{an}})&\text{---}\widetilde{\mathfrak g}=\operatorname{Lie}(\Pi^{\mathrm{alg}}).
\end{aligned}
\]

% block: upper third, lines 6-7 | text
Fonctorialité pour $\Pi$ variable comme pro-objet de la catégorie des groupes profinis.

% block: middle third, lines 8-10 | text
Attention, $\Pi^{\mathrm{an}}$ et \deleted{\unsure{de}} $\pi'^{\mathrm{an}}$ ne sont pas pareils, si $\pi'=\varprojlim\pi_\lambda$. \deleted{\unsure{pro-objet|pour}}\inserted{\unsure{D°|D'} pour $\Pi^{\mathrm{alg}}$, et les $\mathfrak g$, $\widetilde{\mathfrak g}$.} \deleted{\unsure{supposé|sauf} \unsure{par|pour} les morphismes de \unsure{transitivité|transition}}
% block: left margin beside Attention | drawing
\margin{\figure{\begin{tikzpicture}[x=1mm,y=1mm]
\draw (0,17) .. controls (11,20) and (21,17) .. (17,13)
.. controls (12,8) and (0,2) .. (3,0)
.. controls (6,-1) and (11,0) .. (13,1);
\end{tikzpicture}}}

% block: middle third, lines 11-16 | text
\uline{Exemple} 6.10.1. Les $\pi_\lambda$ finis, morphismes de transition surjectifs, $\pi'=\varprojlim\pi_\lambda$ groupe analytique de dim $>0$. Alors $\Pi^{\mathrm{an}}$ n'est autre que $(\pi_\lambda)$ lui-$=$, est formé de groupes discrets, mais $\pi^{\mathrm{an}}$ contient \inserted{\unsure{``}comme facteur''} $\pi'$! [\inserted{Ici} $\Pi^{\mathrm{an}}$ \deleted{plus} est ``plus petit'' que $\pi'^{\mathrm{an}}$!].

% block: lower third, lines 17-23 | text
\uline{Exemple} 6.10.2. Les $\pi_\lambda$ les sous-groupes d'indice fini d'un $\pi$. Alors $\pi'=\{e\}$, donc $\pi'^{\mathrm{an}}$ et $\pi'^{\mathrm{alg}}$ sont les pro-groupes unités. Mais $\Pi^{\mathrm{an}}$ est le système des ss-groupes \deleted{d'indice fini}\inserted{ouverts} des $G_\alpha$ définissant $\pi^{\mathrm{an}}$, et $\Pi^{\mathrm{alg}}$ n'est autre que $\widetilde G^0$, où $\widetilde G=\pi^{\mathrm{alg}}$! \deleted{Donc} [Ici $\Pi^{\mathrm{an}}$ est ``plus grand'' que $\pi'^{\mathrm{an}}$, qui est trivial!]
